%%%%%%%%%%%%%%%%%%%%%%%%%%%%%%%%%%%%%%%%%%%%%%%%%%%%%%%%%%%%%%%%%%%%%%%%%%%%%%%%%%
\begin{frame}[fragile]\frametitle{}
\begin{center}
{\Large Tkinter}

\tiny{(Reference: CS2316, Fall 2016 http://cs2316.gatech.edu/slides/tkinter.pptx)}

\end{center}
\end{frame}



%%%%%%%%%%%%%%%%%%%%%%%%%%%%%%%%%%%%%%%%%%%%%%%%%%%
\begin{frame}[fragile] \frametitle{GUI: Graphical User Interface}

\adjustbox{valign=t}{
\begin{minipage}{0.45\linewidth}
\begin{itemize}
\item A way for people to interact with computer programs
\item Uses windows, icons and menus
\item Can be manipulated by a mouse and keyboard
\end{itemize}
\end{minipage}
}
\hfill
\adjustbox{valign=t}{
\begin{minipage}{0.45\linewidth}
\begin{center}
\includegraphics[width=\linewidth,keepaspectratio]{tkintergui}
\end{center}
\end{minipage}
}

\end{frame}

%%%%%%%%%%%%%%%%%%%%%%%%%%%%%%%%%%%%%%%%%%%%%%%%%%%
\begin{frame}[fragile] \frametitle{GUIs: Step by Step}
\begin{itemize}
\item Design the GUI
\item Create all the buttons, labels, entries, etc. (widgets) And specify where they go in a window
\item Write functions/methods to do what needs to be done
\item Connect the widgets to the functions/methods
\item Start the main loop
\end{itemize}
\end{frame}

%%%%%%%%%%%%%%%%%%%%%%%%%%%%%%%%%%%%%%%%%%%%%%%%%%%
\begin{frame}[fragile] \frametitle{Tkinter}
\begin{itemize}
\item Tkinter is the de facto Python standard GUI library, it comes with your installation of Python and doesn't need to be installed separately
\item We will use the Tkinter GUI toolkit to create Python GUIs \\
https://docs.python.org/3/library/tk.html 
\end{itemize}
\begin{lstlisting}
from tkinter import *
\end{lstlisting}
\begin{block}{Intuition}
\texttt{win.mainloop()} hands control over to Tkinter's event loop -- your script effectively pauses there, waiting and reacting to clicks/keystrokes as they happen, instead of running top-to-bottom like a normal script. This is why widgets are created and \texttt{.pack()}'d \emph{before} \texttt{mainloop()}: nothing shows up on screen until the loop is running.
\end{block}
\end{frame}


%%%%%%%%%%%%%%%%%%%%%%%%%%%%%%%%%%%%%%%%%%%%%%%%%%%
\begin{frame}[fragile] \frametitle{Some simple tkinter widgets}
This code creates a basic GUI with a Label and a Button and places them in a window using the Pack layout manager. The button has no ``command'' so it has no function when clicked.

\adjustbox{valign=t}{
\begin{minipage}{0.45\linewidth}
\begin{itemize}
\item Label: Just displays text, not clickable
\item Button: Can be linked to a function/method which is called on click
\item Entry: Can allow the user to enter text, Can also be used to display text, Has 3 possible states: normal, readonly, disabled
\end{itemize}
\end{minipage}
}
\hfill
\adjustbox{valign=t}{
\begin{minipage}{0.45\linewidth}
\begin{lstlisting}
win = Tk()
win.title("My GUI")
l = Label(win, text="Hello, world")
b = Button(win, text="Click me!")
l.pack()
b.pack()

win.mainloop()
\end{lstlisting}
\begin{center}
\includegraphics[width=0.5\linewidth,keepaspectratio]{tkintersimple}
\end{center}
\end{minipage}
}

\end{frame}


%%%%%%%%%%%%%%%%%%%%%%%%%%%%%%%%%%%%%%%%%%%%%%%%%%%
\begin{frame}[fragile] \frametitle{Some simple tkinter widgets}

\adjustbox{valign=t}{
\begin{minipage}{0.45\linewidth}
\begin{itemize}
\item Now the clicked function will be called every time the button is clicked
\item If we defined the clicked function below the GUI code, the line creating the Button would throw a NameError
\item We can avoid this problem by writing a class for our GUI
\end{itemize}
\end{minipage}
}
\hfill
\adjustbox{valign=t}{
\begin{minipage}{0.45\linewidth}
\begin{lstlisting}
def clicked():
    print("Button clicked!")

win = Tk()
win.title("My GUI")
l = Label(win, text="Hello, world")
b = Button(win, text="Click me!", command=clicked)
l.pack()
b.pack()

win.mainloop()
\end{lstlisting}
\end{minipage}
}
\end{frame}

%%%%%%%%%%%%%%%%%%%%%%%%%%%%%%%%%%%%%%%%%%%%%%%%%%%
\begin{frame}[fragile] \frametitle{Encapsulating a GUI into an Object}

\adjustbox{valign=t}{
\begin{minipage}{0.4\linewidth}
\begin{itemize}
\item The \_\_init\_\_ method takes in a Tk window, creates GUI widgets, and packs them
\item Command is now self.clicked because clicked is a method in the class
\end{itemize}
\end{minipage}
}
\hfill
\adjustbox{valign=t}{
\begin{minipage}{0.5\linewidth}
\begin{lstlisting}
class MyGUI:

	def __init__(self, win):
		self.l = Label(win, text="Hello, world!")
		self.b = Button(win, text="Click me!", command=self.clicked)
		self.l.pack()
		self.b.pack()

	def clicked(self):
		print("Button clicked!")
\end{lstlisting}
\end{minipage}
}
\end{frame}

%%%%%%%%%%%%%%%%%%%%%%%%%%%%%%%%%%%%%%%%%%%%%%%%%%%
\begin{frame}[fragile] \frametitle{Encapsulating a GUI into an Object}

\adjustbox{valign=t}{
\begin{minipage}{0.45\linewidth}
\begin{itemize}
\item In the main method, we create a Tk window, give it a title, and pass it as the win argument of a new GUI object
\end{itemize}
\end{minipage}
}
\hfill
\adjustbox{valign=t}{
\begin{minipage}{0.45\linewidth}
\begin{lstlisting}
def main(args):
    win = Tk()
    win.title("My GUI")
    MyGUI(win)
    win.mainloop()

if __name__ == '__main__':
    import sys
    main(sys.argv)
\end{lstlisting}
\end{minipage}
}
\end{frame}

%%%%%%%%%%%%%%%%%%%%%%%%%%%%%%%%%%%%%%%%%%%%%%%%%%%
\begin{frame}[fragile] \frametitle{GUI layout managers}

\adjustbox{valign=t}{
\begin{minipage}{0.45\linewidth}
Pack
\begin{itemize}
\item Positions widgets relative to each other
\item Simple to use
\item Limited layout possibilities compared to grid
\end{itemize}
\end{minipage}
}
\hfill
\adjustbox{valign=t}{
\begin{minipage}{0.45\linewidth}
Grid
\begin{itemize}
\item Allows you to place widgets in rows and columns
\item More complicated than pack but good for GUIs where a particular arrangement is desired
\end{itemize}
\end{minipage}
}

Don't use grid and pack in the same container!
\end{frame}

%%%%%%%%%%%%%%%%%%%%%%%%%%%%%%%%%%%%%%%%%%%%%%%%%%%
\begin{frame}[fragile] \frametitle{Exercise: Adder}

\begin{itemize}
\item Create a GUI with two NORMAL entries and one READONLY entry
\item There should a button that, when clicked, adds the numbers entered in the first two entries and displays the sum in the third entry
\item If the calculation can't be performed due to invalid input, display a message box with an error message.

\end{itemize}

\begin{center}
\includegraphics[width=\linewidth,keepaspectratio]{adder}
\end{center}
\end{frame}

%%%%%%%%%%%%%%%%%%%%%%%%%%%%%%%%%%%%%%%%%%%%%%%%%%%
\begin{frame}[fragile] \frametitle{Frames}

\begin{itemize}
\item Container that can be placed within a window
\item Holds other widgets
\item Each frame has its own grid layout, independent of the main window
\item You can use a different layout manager for each container
\end{itemize}
\end{frame}

%%%%%%%%%%%%%%%%%%%%%%%%%%%%%%%%%%%%%%%%%%%%%%%%%%%
\begin{frame}[fragile] \frametitle{Example: Frames (use within GUI class)}
\begin{lstlisting}
f1 = Frame(win, height=100, width=100, bd=1, relief=SUNKEN)
f2 = Frame(win, height=100, width=100, bd=1, relief=SUNKEN)

f1.grid(row=0, column=0, sticky=EW)
f2.grid(row=0, column=1, sticky=EW)

# self.l1 will be used for another method later
self.l1 = Label(f1, text="One", height=5, width=5)
self.l1.grid(row=0, column=0)
Label(f1, text="Two", height=5, width=5).grid(row=0, column=1)
Label(f2, text="Three", height=5, width=5).pack(side=LEFT)
Label(f2, text="Four", height=5, width=5).pack(side=LEFT)

\end{lstlisting}
\end{frame}

%%%%%%%%%%%%%%%%%%%%%%%%%%%%%%%%%%%%%%%%%%%%%%%%%%%
\begin{frame}[fragile] \frametitle{Radiobuttons}
\begin{itemize}
\item Allows the user to choose from one or more options
\item Can be linked to other Radiobuttons with StringVars and IntVars
\item If configured correctly, only one Radiobuttons from a group can be selected at at time
\end{itemize}
\end{frame}

%%%%%%%%%%%%%%%%%%%%%%%%%%%%%%%%%%%%%%%%%%%%%%%%%%%
\begin{frame}[fragile] \frametitle{StringVars and IntVars}
\begin{itemize}
\item Two variable classes included in Tkinter
\item Have .set() method to set their value
e.g. my\_string\_var.set(``Hello'')
\item StringVars have string value, IntVars have integer value
\item Have .get() method to get their value
e.g. val = my\_string\_var.get()
\item Can be passed in as ``variable'' argument when creating a Radiobutton, or ``textvariable'' argument when creating an Entry

\end{itemize}
\begin{block}{Intuition}
A plain Python variable is just a value -- assigning to it doesn't touch the screen. A \texttt{StringVar}/\texttt{IntVar} is a small wrapper object that Tkinter watches: link it to a widget via \texttt{variable}/\texttt{textvariable}, and the widget updates automatically whenever \texttt{.set()} changes the wrapper, in either direction.
\end{block}
\end{frame}

%%%%%%%%%%%%%%%%%%%%%%%%%%%%%%%%%%%%%%%%%%%%%%%%%%%
\begin{frame}[fragile] \frametitle{Exercise}
\begin{itemize}
\item Revisit our Adder example from the last class
\item Instead of using Entry.insert and Entry.delete, use a StringVar to change the text of the result Entry
\item See what happens if you don't set the Entry state to ``normal'' before changing the text!
\end{itemize}
\end{frame}

%%%%%%%%%%%%%%%%%%%%%%%%%%%%%%%%%%%%%%%%%%%%%%%%%%%
\begin{frame}[fragile] \frametitle{Adding Radiobuttons to our Frame example}
\begin{lstlisting}
self.color = StringVar()

rb1 = Radiobutton(win, text="Red", variable=self.color, value="red")
rb2 = Radiobutton(win, text="Blue", variable=self.color, value="blue")
rb3 = Radiobutton(win, text="Green", variable=self.color, value="green")

self.color.set("unknown") # deselects all radiobuttons linked to self.color

rb1.grid(row=1, column=0)
rb2.grid(row=2, column=0)
rb3.grid(row=3, column=0)

Button(win, text="Change Color", command=self.change_color).grid(row=1, column=1)
\end{lstlisting}
\end{frame}

%%%%%%%%%%%%%%%%%%%%%%%%%%%%%%%%%%%%%%%%%%%%%%%%%%%
\begin{frame}[fragile] \frametitle{Adding Radiobuttons to our Frame example}

\adjustbox{valign=t}{
\begin{minipage}{0.45\linewidth}
\begin{lstlisting}
def change_color(self):
	new_color = self.color.get()
	if new_color == "unknown":
		pass
	else:
		self.l1.config(bg=new_color)
\end{lstlisting}
\end{minipage}
}
\hfill
\adjustbox{valign=t}{
\begin{minipage}{0.45\linewidth}
\begin{center}
\includegraphics[width=\linewidth,keepaspectratio]{rb1}
\end{center}
\end{minipage}
}
\end{frame}

%%%%%%%%%%%%%%%%%%%%%%%%%%%%%%%%%%%%%%%%%%%%%%%%%%%
\begin{frame}[fragile] \frametitle{Additional Exercise}
\begin{itemize}
\item Add another set of Radiobuttons to the GUI. Separate from those linked to self.change\_color
\item There should be four Radiobuttons, each linked to one of the labels
\item Create an IntVar to use as the variable for these Radiobuttons
\item The user's selection should determine which Label will change color when the button is clicked.
\end{itemize}
\end{frame}

%%%%%%%%%%%%%%%%%%%%%%%%%%%%%%%%%%%%%%%%%%%%%%%%%%%
\begin{frame}[fragile] \frametitle{Root Window vs. Toplevels}
\begin{itemize}
\item The root window is your main window (Tk)
\item A GUI application should have only one root window
\item Any additional windows should be Toplevels
\item Hide windows with .withdraw()
\item Show windows with .deiconify()
\item Delete windows with .destroy() (they will be gone forever!)
\item If you destroy the root window, all Toplevels will also be destroyed

\end{itemize}
\end{frame}

%%%%%%%%%%%%%%%%%%%%%%%%%%%%%%%%%%%%%%%%%%%%%%%%%%%
\begin{frame}[fragile] \frametitle{Adding multiple windows to our example}
\begin{lstlisting}
# in __init__ method, add line: self.root = win

def new_window(self):
	self.root.withdraw()
	self.new = Toplevel()
	Label(self.new, width=30, height=10, text="I'm a new window!").pack()
	Button(self.new, text="Close window", command=self.close_new).pack()

def close_new(self):
	self.new.destroy() # destroy new window
	self.root.deiconify() # show root window

\end{lstlisting}
\end{frame}


%%%%%%%%%%%%%%%%%%%%%%%%%%%%%%%%%%%%%%%%%%%%%%%%%%%
\begin{frame}[fragile] \frametitle{Changing the GUI appearance}
\begin{itemize}
\item GUI windows and widgets have many optional parameters to change their appearance
\item bg: background color
\item fg: foreground color (for Labels, this is the font color)
\item bd: border width
\item border: border type (RAISED, SUNKEN, etc)
\item For color codes: http://htmlcolorcodes.com/
\item For more options, check the documentation or online tutorials
\end{itemize}
\end{frame}

%%%%%%%%%%%%%%%%%%%%%%%%%%%%%%%%%%%%%%%%%%%%%%%%%%%
\begin{frame}[fragile]\frametitle{Quick Check: Tkinter}
\begin{block}{Think About It}
The Adder exercise wraps its GUI in a class (\texttt{MyGUI}) instead of writing the widget-creation code directly at module level, the way the very first Label/Button example did. What breaks if you don't do that once a callback like \texttt{clicked()} needs to read or update \emph{other} widgets?
\end{block}
\pause
\begin{block}{Answer}
A bare function callback only has access to whatever it can see via globals or closures -- referencing widgets created later in the same script raises a \texttt{NameError}, and referencing several widgets from many callbacks gets messy fast. Wrapping everything in a class turns each widget into \texttt{self.widget}, so any method (bound via \texttt{command=self.clicked}) can freely read and update every other widget through \texttt{self}, regardless of creation order.
\end{block}
\end{frame}

